% =====================================================================
%  R. Nielsen, IDEA AND REALITY
%  In: For Idea and Reality.  A Journal, edited by R. Nielsen, Bj. Bjørnson and Rud. Schmidt.  Volume for 1869.
%  Copenhagen: Chr. Steen & Søn; printed by Louis Klein.  pp. 1-39.  (Ide og Virkelighed, in For Ide og Virkelighed, Aargang 1869.)
%  ENGLISH TRANSLATION, from transcription.tex (Danish).
%  GENERATED by ebook-method/for-idee-1869/tr_build.py --write from this template and the translators' parts in
%  ebook-method/for-idee-1869/tr/ (TRANSLATOR-BRIEF.md; terms: GLOSSARY.md, the titles and four key terms approved by
%  Hans 2026-10-13).  Conventions: TRANSLATION-PLAYBOOK.md.  Page numbers of the 1869 print in the margin.
% =====================================================================
\documentclass[12pt,a4paper,openany]{book}
\usepackage[utf8]{inputenc}
\usepackage[T1]{fontenc}
\usepackage{libertinus}
\usepackage{libertinust1math}
\usepackage[english]{babel}
\usepackage[protrusion=true,expansion=true]{microtype}
\usepackage{setspace}
\usepackage[margin=1.3in]{geometry}
\usepackage{fancyhdr}
\usepackage{hyperref}
\hypersetup{pdftitle={Idea and Reality}, pdfauthor={Rasmus Nielsen}, hidelinks}
\setlength{\headheight}{15pt}
\pagestyle{fancy}
\fancyhf{}
\fancyhead[LE,RO]{\thepage}
\fancyhead[C]{\textit{R. Nielsen: Idea and Reality (1869)}}
\setstretch{1.3}
\usepackage{marginnote}
\renewcommand{\thefootnote}{\arabic{footnote})}   % the print numbers each page's note ¹) (p. 553: *))
\newcommand{\opage}[1]{\marginnote{\footnotesize\textit{[#1]}}}
\newcommand{\essayhead}[2]{\begin{center}{\large #1}\ifx&#2&\else\\[0.3em]{\normalsize #2}\fi\\[0.5em]\rule{3em}{0.4pt}\end{center}\medskip\noindent}
\newcommand{\signature}[1]{\par\smallskip\hfill #1\hspace*{2em}\par}

\begin{document}

% the volume's title page (PDF 7)
\begin{titlepage}
\begin{center}
\vspace*{3em}
{\Huge For Idea and Reality.}\\[1em]
\rule{8em}{0.4pt}\\[1em]
{\large A Journal,}\\[0.8em]
{\small edited}\\[0.4em]
{\footnotesize by}\\[0.6em]
{\large R. Nielsen,}\\[0.4em]
{\large Bj. Bjørnson and Rud. Schmidt.}\\[1em]
\rule{8em}{0.4pt}\\[1em]
{\large Volume for 1869.}\\[1em]
\rule{8em}{0.4pt}\\[4em]
Copenhagen.\\[0.4em]
Chr. Steen \& Søn, Publishers.\\[1.5em]
{\small Printed by Louis Klein.}\\[3em]
{\small Translated from the Danish}
\end{center}
\end{titlepage}

\essayhead{Idea and Reality.}{}
% --- p. 1 ---
\setcounter{footnote}{0}%
\opage{1}Idea and reality are a peculiar matter: they cannot
do without one another and yet are constantly at strife
with one another. Reality by itself, without guiding ideas
and luminous principal thoughts, presents only a dismal
spectacle, a ghastly sight of divided interests, petty
squabbling, confusion and dissolution; and yet under such
circumstances it is by no means convenient for peevish reality
to engage with ideas, least of all with
the idea it most deeply lacks, the awakening, clarifying, rejuvenating
idea, for which, without understanding itself,
it has long sighed in secret.

Reality's protest against a new, advancing
idea is no valid proof of the idea's invalidity.
On the contrary! A side of life may in and for itself be perfectly
justified, may have every title as a pretender
to the throne of reality, and yet meet resistance
in the circumstances and be repeatedly
beaten back, either because practical life still
stands under the dominion of antiquated traditions, venerable
by their age, which must die away only gradually,
or else because reality, by suddenly
shaking off a whole dynasty of old, forfeited
% --- p. 2 ---
\setcounter{footnote}{0}%
\opage{2}ideas, has had its blood set boiling and has fallen into such
an uproar that some time must pass before it
can properly collect itself and provide the conditions for the
new, vigorous idea. At last the long-awaited
hour comes, the happy moment of understanding, reality
opens its eye and recognizes the idea: idea and
reality rush into one another's arms like two lovers
who have long had to struggle with misunderstandings. But
happy moments! alas, they were only brief. When
the oath of fidelity is sworn, the festivities of the consecration over,
and the honeymoon at an end, then enthusiasm cools,
the creditors of existence present their bill
for the conditions, and the troubles come again;
for that is, after all, reality's lot: it is always
full of troubles. Now complaint follows upon complaint from both
sides. The idea complains that reality is sluggish, trivial % print: Virkelighen (as printed)
and spiritless, without earnestness, without perseverance, without uplift.
Reality complains that the idea, with its far-reaching
demands and high-flying plans, is impractical,
reckless, with no eye for what lies nearest, for the useful, an
idealist, a visionary, in plain words, a dreamer. There may
be violent scenes between the pair; the strife may go
so far that, despite mutual obligations and
the solemn vows of marriage, they even appeal to the public's
most esteemed sense of justice and demand a divorce.

To judge in the case, it is each time necessary
to investigate on which side the guilt lies.
But to find out on which side the guilt lies, one must
in every case concerned be able to distinguish
between a true and a false idea. If it is a false idea
that has played the seducer's part and by dazzling
% --- p. 3 ---
\setcounter{footnote}{0}%
\opage{3}appearance and flattering phrases has beguiled reality
in an unguarded moment and tricked it onto a dangerous and
ruinous path, then it is plainly a duty to come to
the poor beguiled one's aid, to chastise the seducer, and
to see, as far as possible, that the hasty step is made good
again. But if the idea is within its good right, if the
goal it has set is the only reasonable one, because it
is the only one that truly answers to the need of the age,
and the path it has pointed out the only right one, because it
is the only one that leads to the goal: then reality is surely
in the wrong when, because it costs effort,
it balks at applying its powers and working undaunted
at the solution of the tasks the idea has prescribed.

What this means may be seen by a glance at the
idea that the present age itself has exalted to be the
ruling one, namely the idea of freedom. Is freedom,
political freedom of the people, when we examine the matter impartially
and speak without circumlocution, a true, self-valid and justified
idea, or perhaps false, self-contradictory and ruinous?
To make this an open question is, in view
of the difficulties, troubles and complications
with which it has already burdened reality, surely not
so untimely as it might seem at first glance.
Yet it is of no use to begin with reality,
for reality is subject to many influences;
reality has many minds, puts on many
garbs and speaks with many tongues. How in the year
1848 the whole real contemporary world was enthusiastic for freedom,
ready to bring every sacrifice and to risk everything
for the idea! That was the time of awakening, the golden time of
% --- p. 4 ---
\setcounter{footnote}{0}%
\opage{4}expectations, the time of hope, and hope was like the
fresh spring, hope was light green. But since that time expectations
have sunk considerably again.
Reality has already several times changed
color; the light green has at times faded into
the grayish and has even in some places passed into
the flame-yellow. No, from reality one gets no
reliable explanation; one must begin with the idea.

But is there then really in our days any reasonable
person who can seriously think of casting
doubt on the political validity of the freedom of the people? We answer
Yes! There are still quite a few reasonable people
who harbor some doubt. In proof of the correctness
of our assertion we shall adduce an original
doubter in this respect, a thinker of the first rank, a clear-sighted, liberal-minded
man, a man who has the unanimous testimony
of his contemporaries that he at least has kept himself outside
all political parties and to that extent must be called disinterested,
a man of superior shrewdness and with a sure eye for
the troubles of reality, a man who precisely set
as his life's task to think through the highest idea of freedom,
and then, remarkably enough, came --- as it seems
--- to the result that freedom of the people, politically understood,
is a selfish delusion, a false idea. We mean S.
Kierkegaard.

S. Kierkegaard was, as is well known, no admirer
of the demands of the age, no admirer of the so-called liberal
endeavors. Already his first articles on ``The Morning Observations
of the Kjøbenhavnspost'' and ``The Polemic
of Fædrelandet'' together with the replies he exchanged on that occasion % print: Polemik» (as printed)
with ``Mr. Orla Lehmann'' are characteristic. It was a
% --- p. 5 ---
\setcounter{footnote}{0}%
\opage{5}polemic of mood, in which he began by trying
his sharp pen; it was his first grip on the string-play
of ideality, a few quick strokes of the bow, to hear the tone,
and the tone indicated plainly enough that this violin would not
tune itself to the blowing instruments of politics.

At a later stage, where the discussion is especially of
state and church, he expresses himself, occasioned by
a remark of Dr. Rudelbach, very sharply about what
he calls an unhappy confusion of politics and Christianity,
the inverted Reformation, which reformingly puts a
new worse in place of an old better, ``though
it is supposed to be certain and true that it is a Reformation
on the occasion of which the whole town is illuminated.'' The
true freedom is not an outward arrangement but
an inward matter; freedom is essentially inwardness. ``As
it is man's advantage over the animal to be able to live
under any climate, so perfection
in Christianity ranks, precisely because it is inwardness, according to its
strength to be able to live, if need be, under the most imperfect
forms. Politics is this outwardness,
this Tantalean busyness for change in the outward.''

The cognition of freedom is in the strictest sense self-knowledge;
liberation can therefore not begin with
anything outward, not take its point of departure either in
the race or in world history, neither in new constitutions of state nor
of church; it must begin from within,
each single one beginning with himself, not with
great expectations of ``the free institutions'' , but in
faith in the liberating spirit of truth.

It is a delusion of the senses when anyone thinks that it
is the free institutions that are to save the Church and
% --- p. 6 ---
\setcounter{footnote}{0}%
\opage{6}Christianity. ``Should this'', says K., ``belong to
true Christianity, this faith in the saving power of
the free institutions politically obtained, then I am no
Christian, indeed worse still, I am a pure child of the Devil;
for to speak honestly, I even suspect these politically
obtained free institutions, especially as to their saving
and regenerating power. That is my Christianity,
or so Christian-ish am I, who by the way has never
concerned myself with `Church' and `State' --- that is far
too great for me; for that quite other
seers than I are required, or else it ought wholly to be left to
those duly appointed and called to it. I
have never fought for the emancipation of `the Church', any more
than for that of the Greenland trade, of women, of the Jews,
or for any other emancipation whatever. I have only as
a single one, with the aim `the single individual', aiming at inwardizing
Christianity in `the single individual', consistently with
the weapons of the spirit, solely and alone with the weapons of the spirit,
fought to draw attention to, and against letting
oneself be deceived by, delusions of the senses.''

As it is then a delusion of the senses if anyone thinks
that it is the outward forms that hinder him from attaining
freedom, so it is also a delusion of the senses, indeed the
same turned in another way, if anyone imagines
that it is the outward forms that are to help him
to freedom.

According to this view S. K. could then neither
attach great hopes to the free movements that
distinguished the year 1848. His little writing: ``On My
Work as an Author'' (1851) ends with the following passage, remarkable for his whole
standpoint: ``They snapped, the
% --- p. 7 ---
\setcounter{footnote}{0}%
\opage{7}year '48, the threads of prudence; the whine was heard
that heralds Chaos! `That was the year '48, that was
progress', yes --- if, that is, government is procured,
for which perhaps not a single
new official is needed, nor the dismissal of any older one, but
perhaps indeed an inward transformation in the direction of becoming
steadfast with God. The fault from above was presumably this, that on the
whole, from the top to the bottom of the government,
the strength was essentially worldly prudence, which precisely
is essentially a lack of strength; the guilt from below
was this, to want to have all government away: the punishment ---
for the way of sin is always the way of punishment --- the punishment
is: what is most bitterly missed is precisely government.
Never as in our century were the race and the
single individuals in it (the commanding --- the obeying; the superior
--- the subordinate; the teacher --- the learner,
etc.) so freed from, if you will, all
encumbrances of there being something that stands and is to stand unconditionally fast;
never did `opinions' (the most heterogeneous,
in the most various domains) feel themselves so in `Liberty, Equality
and Fraternity', unembarrassed and made blissful under
the free pass of `to a certain degree'; never will the race so
deeply come to see the truth that it and each single one in it
needs something to stand unconditionally fast, needs
what the Deity, the loving and in love,
devised, the unconditional, in place of which man,
the clever to his own corruption, admiring himself, set
this admired `to a certain degree'. Bid the sailor
sail without ballast --- he capsizes; let the race, let
each single one in it try to subsist without the unconditional:
it is and remains a whirl. It may for a longer or
% --- p. 8 ---
\setcounter{footnote}{0}%
\opage{8}shorter time, at intervals, seem otherwise, that it is
firmness and security, at bottom it is and remains a whirl,
even the greatest events and the most strenuous
life nevertheless a whirl, or like sewing without fastening the
end --- until the end is again made fast by this, that the
unconditional is introduced, or by this, that the single individual, even
at ever so great a distance, yet relates himself to an
unconditional. To live merely in the unconditional, to breathe only
the unconditional, man cannot, he perishes
like the fish when it is to live in the air; but on the
other hand, without relating himself to the unconditional a
man cannot in a deeper sense `live' either, he
expires, that is, he perhaps continues to live,
but spiritlessly. If --- to stay with what is mine: the religious
--- the race or a great part of single individuals in the race
have outgrown this childlike thing, that another man for
them is the one who represents the unconditional: well,
the unconditional can nevertheless not be dispensed with. Then
the single individuals must relate themselves to the unconditional themselves.
It is this that I, according to my granted ability, with utmost exertion
and many a sacrifice, have fought for, fighting
against every tyranny, also that of the numerical.
This striving of mine has been made an object of hatred as immense
pride and arrogance --- I believed, and I believe, that this
is Christianity and love of the neighbor.''

However a free constitution may be arranged --- here we speak
of free constitutions in all generality --- it must in the end
go out over both state and citizens. The authority of government
is weakened. No single one either dares or can
step forward decisively, act from his own conviction and
take the responsibility upon himself; but everything is conditioned by the
% --- p. 9 ---
\setcounter{footnote}{0}%
\opage{9}mood of the moment, everything is made dependent on the consent of the majority.
The damage is twofold. In the first place the expert must
put up with it that his proper principal thought, his carefully
worked-out plan is clipped and pruned under oppositions and discussions,
and by obstinate objections,
partial alterations, accommodating adaptations
is so botched that, far from becoming of one casting,
it becomes on the contrary a patchwork of heterogeneous pieces,
a motley composition of unlike patterns, less
like a plan than a quilt. When such a patchwork
has been put through, and bitter experience has at long last
shown how useless it is: whose then is the
guilt? No one's! Who bears the responsibility? No one! This
is still the very worst. Responsibility evaporates in
discussions, vanishes in committee reports and votes,
so that the whole comes out at last as
the work of a machine, a kind of anonymous product of the
driving wheels of the parties. If one tries to put a better in
the place of the recognized worse, then the
political comedy begins afresh. Everything goes on publicly and is
yet a continued game of intrigue; it looks as if each
single matter were subjected to an all-sided trial, a careful, extensive
grounding; but what is it that
is actually tried? It is not the thing itself, but the parties'
forces. What gives the grounds weight? Not the nature
of the matter, not convincing thoughts, but the voters'
wholly or half masked tendencies. What is in and
for itself true, reasonable, useful, naturally cannot
without further ado come into consideration; for the conditions are
far more important than the unconditional. It is simply not
done to speak out, to defend it, according to one's
% --- p. 10 ---
\setcounter{footnote}{0}%
\opage{10}own conviction, as the only right thing, leaving it to the surroundings
what they will accept or reject, no, the point
is to catch the wind, to stretch out the feelers, to sound out
the parties, to learn what one can put through. To get something put through is the most important; the
next most important is what it is: half a truth,
if only it can go through, is better than a whole one
that exposes itself to being rejected. In this way
everything becomes acceptable and everything rejectable to a certain degree,
truth to a certain degree, and consequently also
untruth, justice to a certain degree, consequently also
injustice, honesty to a certain degree, consequently also
dishonesty. To be straightforwardly untrue, unjust, dishonest
will not do; but to be unconditionally true, unconditionally
just, open and honest will not do either: reality
can use pure truth as little as pure
untruth; it demands a peculiar mixture of good
and bad, of honesty and dishonesty. The character of reality
becomes in this sense a shrewd experience,
a cleverly calculated characterlessness. But the clever characterlessness
is moral ambiguity, a refined subtlety
that does not attack openly but slyly
hollows out what exists. This is for the civic way of thinking
on the whole the lulling, stupefying, demoralizing thing.
``For against rebellion one can use force,
for notorious false speech the punishment is in store; but dialectical
secretiveness is hard to root out;
a comparatively finer ear is needed to trace
the soundless sneaking of reflection along the byways of ambiguity.
What exists exists, but the passionlessness of reflection
finds its reassurance in this, that it is an ambiguity.
% --- p. 11 ---
\setcounter{footnote}{0}%
\opage{11}One does not want the monarchy abolished, by no means,
but if one could little by little get it transformed
into an imagination, then one would gladly cry
Hurrah for the King. One does not want distinction
overthrown, but if at the same time one could promote
a knowledge that it is an imagination --- then one would
admire. One wants to let the whole Christian terminology
stand, but secretly to be aware
that nothing decisive is to be thought by it. And
one wants to be without regret; for one tears down nothing.''

To the darkest shadow-regions of free constitutions belongs
precisely this, that the will of the people must always, mediately or immediately,
make itself known as the will of the crowd, and
the crowd is, this is K.'s thesis, ``The crowd is untruth.''
At bottom there are none who have more contempt for man
as man than those ``who make a profession of standing
at the head of the crowd. Let one come to such a one,
a single human being --- well indeed, what does he care
for that; it is much too little; proudly he turns him away;
there must be a hundred at least. And when there are
a thousand, then he bows before the crowd, bows and scrapes;
what untruth! --- That the crowd as crowd is
all, the single individual in the crowd nothing, means that the crowd
is the public. But precisely the circumstance that the
single individual in the crowd is nothing makes it possible that he
can identify himself with the crowd, speak and write in the name of the
crowd without standing personally to account for
what he says. For who is it one has before one?
The crowd and yet no one, an anonymous. ``That an anonymous one,
by help of the press, day in and day out can get
(also in respect of the intellectual, the ethical,
% --- p. 12 ---
\setcounter{footnote}{0}%
\opage{12}the religious) said what he wants, what he perhaps
would not in the remotest way have had the courage to say even the
least of in the situation of singleness; can every time
he opens his --- well, mouth one cannot call it
--- his gullet, \emph{at one stroke} address himself to a thousand
times a thousand; can get a thousand times a thousand
to say after him what was said --- and no one has the responsibility,''
this K. regards as the great misery of our age, a misery
that is greater than the greatest of antiquity. ``Let no one say
that truth can after all by help of the press
overtake the lie and the error again. You who speak
thus, ask yourself: do you dare to maintain that men
taken as a crowd are just as quick to reach
for truth, which is not savory, as for
untruth, which is always daintily prepared, to say nothing of when
it is moreover to be bound up with an admission
that one has let oneself be deceived! Or do you merely dare
to maintain that truth is as quickly understood
as untruth, which demands no prior knowledge,
no school, no discipline, no abstinence, no self-denial, no honest self-concern, no slow
labor! No, truth, which also abhors this
untruth: as sole aim to seek diffusion,
is not so quick on its feet.'' (``The Point of View for My
Work as an Author.'' pp. 95--96.)

The dark picture is completed in the depiction of the public.
Certainly one can by no means say that it is the idea
of the freedom of the people that produces the public; for that is culture:
the public of the present age is a peculiar degeneration
of the culture of the present age. On the other hand one can say with good reason
that the free constitution has its most ambiguous friend, its most
% --- p. 13 ---
\setcounter{footnote}{0}%
\opage{13}dangerous ally in the public. What then is the public?
J. L. Heiberg, who, as is well known, has fought several
hot battles with the public, has also delivered valuable
contributions to its description. In his treatise ``People and
Public'' (Intelligensblade, 1842, No. 6) he proceeds
from the well-known opposition, that a people is an organism,
a public a mass, and expresses himself on the transition
from people to public thus: ``The successive
transformation of the people into a public is a fact which in this country
already has its \emph{history}, and I myself have delivered
not a few documents toward it, though partly without
being aware of it, for since our political development is
much younger than our literary one, but it is precisely the latter
that first throws the right light on the former, one could
long be ignorant of the true significance of
the forces moving in literature. Already in
``Christmas Jests and New Year's Pranks'' I introduced the public personified
as ``a colossal monstrosity'', but let
the poet in the first moment greet it with the words: ``My
people!'' Only later, when at the sight of its thousands
of eyes, mouths and hands, and by its crude, self-contradicting
utterances he gets sight of its atomistic
nature, he meets it with the utmost contempt,
without however taking back the expression with which he first
greeted it.''

It is now easily understood that public opinion, which
plays so great a role in all free commonwealths, can
under the semblance that it is the voice of the people be appropriated
by the public. Precisely because the public is always spectator and
the single individual in the public essentially nothing --- insofar
as he is otherwise something, perhaps even something capable,
% --- p. 14 ---
\setcounter{footnote}{0}%
\opage{14}steps at once outside the public --- then the public is
not acting; its only action is to make its
opinion known in expressions of approval and displeasure. The making-known
itself is inarticulate. It is a chattering
and prattling for and against, a humming and droning, a whirring
and buzzing, to which a politician must prick up his ears
every time he wants to listen out whether it is approval or displeasure
that actually has the preponderance. What conjures up
the phantom of a public opinion is no whole, simple
thought, but a magnitude of the utmost complexity, a swarm
of countless little infusorial opinions, a sum of fractions
with different denominators and different signs.
That such a collective magnitude can impossibly yield
a view in which there is popular unity, still less
be determinative of a true will of the people, is obvious.
``Only in a passionless but reflective age'' --- says
S. Kierkegaard, who in this respect honestly continued
what J. L. Heiberg had begun --- ``can that phantom
which is called the public develop by means of the press,
when the latter itself becomes an abstraction. In enthusiastic
times, in passionate tumultuous times, even when a
people wishes to realize the idea of the barren desert, to raze
everything to the ground: still there is no public\dots{} The public is a
concept which cannot occur at all in antiquity, because
the people itself, en masse, in corpore, had to step forward into
the situation of action, had to bear the responsibility for what
the single individual from its midst brought about, while again the
single individual --- personally, as this determinate person --- had to be
present, had to submit to the court of the moment in
approval or disapproval.''

That the public cannot be the people itself, indeed cannot even
% --- p. 15 ---
\setcounter{footnote}{0}%
\opage{15}represent the people, is evident, among other things, from the fact
that the public cannot act. The moment the public
wants to come forward acting, it ceases to be the public;
one may then, as circumstances dictate, call it mass,
crowd, rabble, people, whatever one likes, but the public it
is not. Where there is talk of a revolution it
is therefore never said that the public has been fired upon, but
that the mass has been fired upon, the people: a phantom
one cannot fire upon. Nor will history ever
be able to report that a public has either
overthrown a throne or driven a prince into exile;
one might as well suppose that it was ghostly mischief that
had overthrown the throne, that a ghost had driven out
the prince. In relation to everything that belongs to action
and character, the public is a ghost. The public's
character is reflecting, ratiocinating idleness,
and its changing opinions are idle. Yes, so ghostlike
idle are these opinions that it is even idle, as S. K.
has aptly shown, to want to be of the same
opinion as the public. ``To adopt the same opinion as
such and such determinate human beings means that one
knows that these will be exposed to the same dangers as
oneself, that they will go astray along with one, if the opinion
is misleading, and so on. But to adopt the same opinion
as the public is a treacherous consolation, for the public is
only in abstracto. Whereas therefore no majority
has ever been so sure of keeping right
and victory as the public is, this is but little
comforting to the single individual, for the public is a phantom
which permits no personal approach. If
% --- p. 16 ---
\setcounter{footnote}{0}%
\opage{16}one today adopts the opinion of the public and tomorrow is
hissed off, then he is hissed off by the public.''

The public can organize nothing, it can only level:
levelling is ``the envy that is establishing itself.''
Levelling is not the act of a single person, but a play of reflection
in the hand of an abstract power; this abstract
power is the public. ``The abstraction of the press (for
a newspaper, a daily, is no civic concretion and
only in an abstract sense an individual) begets the abstraction's
phantom: the public, which is the properly levelling thing\dots{}.
The less idea there is in an age, the more
it, exhausted by a flaring-up enthusiasm, takes its rest
in indolence, even if we were to imagine that the press became
weaker and weaker because no event, no idea
grips the age: the more easily will levelling become
a pernicious pleasure, a titillation of the senses, which for a moment tickles
and only makes evil worse. And if one has often
described the demoralization of absolute monarchy, the decay
of revolutionary times, then the decay of a passionless age
is something just as pernicious, even if, owing to its ambiguity,
less conspicuous. This sluggish crowd,
which understands nothing itself and wants to do nothing itself, this
gallery-public now seeks pastime and then gives itself up
to the conceit that everything anyone does is done
so that it may get something to talk about. Sluggishness
sits grandly with its legs crossed, and everyone who wants to
work, the king and the official and the people's teacher
and the abler journalist and the poet and the artist,
all become as if harnessed in order to drag this
sluggishness forward, which grandly believes that the others are the horses.
--- So the public keeps itself a dog for
% --- p. 17 ---
\setcounter{footnote}{0}%
\opage{17}amusement's sake. This dog is literary contemptibility.
Now if a better man shows himself, perhaps even an excellent one,
the dog is set on him, and the fun begins.
The snappish dog tears at the skirts of his coat, allows itself
every unruly impertinence --- until the public grows tired
of it and says: now that may be enough. Then the public has
levelled. The better man, the stronger man has been maltreated, and
the dog, well, it goes on being a dog, which the public
itself despises\dots{}. And the public is without remorse, for it
is not the public --- it was the dog; just as
one says to children: it was the cat. And the public is
without remorse, for it was no disparagement --- it was
merely a little fun. And the public wants to be without remorse, for
it is really not they who keep the dog, they merely subscribe,
nor do they set it on directly, nor
do they whistle at it outright; in case of a lawsuit
the public would say: the dog is not mine, it is
a stray dog; and in case the dog were
sent to the veterinary school to be put down, the
public could still say: it was really good that the
nasty dog was killed, we have all wished it --- even
the subscribers.'' (A Literary Review. pp. 95-96.)

The result is easy to draw: the freedom of the people is conditioned
by this, that every unconditional authority is to a certain degree weakened
and broken down, so that all authority, including governmental authority,
sinks little by little beneath the level of relativity.
Through tough but pliant constitutional forms, movable in all
directions, the parties with their leaders, making the best possible use of the oscillations
of public opinion, work to promote the interests
which, whether they are otherwise harmful
% --- p. 18 ---
\setcounter{footnote}{0}%
\opage{18}or beneficial to society, rational or irrational, at
a given time show themselves dominant with the majority. The majority's
opinion is, to be sure, officially stamped as valid
for the opinion of the people, but ought nevertheless, since it each
time must seek its proper center of gravity in the preponderant
majority, that is, in the crowd, most of all to be conceived
as a constitutional, though provisional, expression
of the crowd's opinion. Then the metamorphoses follow as
of themselves: the government becomes a people, the people becomes a
crowd, the crowd becomes a public, the public becomes
a phantom, a political ghost, a kind of indeterminable
nothing. These are retrogressive metamorphoses, premises
from which it is concluded: ergo the freedom of the people is a
vain conceit, a false idea.

Now if this conclusion were wholly reliable
and the result gained in all and everything irrefutable:
then it would certainly look sorry in the present, not
least with us. But the conclusion is, as the attentive
reader has surely already told himself, not reliable.
In the ambiguity of the premises there is a composition
of heterogeneous elements, a dubious
mixture, which can be separated and cleared up only when we reflect
on the relation between idea and reality.

Those who, under uncanny impressions of what
offends them in reality, have grown accustomed to adjust
their judgments to their mood and on that account
think that under the pressure of the freedom of the people they really have
grounds to be dispirited, will find in an author such as
S. Kierkegaard utterances that correspond to their personal
view, soothing flashes of wit, polemical, refreshing
sallies, expressions so apt that it must at times
% --- p. 19 ---
\setcounter{footnote}{0}%
\opage{19}seem to them as if he had read in their political
heart and taken every word out of their own inner being; they
can, so far, read him with a spiritual-aristocratic self-regard
for ironic edification and amuse themselves with the
witty images; --- but to declare him a decided opponent
of the freedom of the people and straightforwardly to claim him as their own
they cannot. With him it is a peculiar matter; in his case
attack essentially means defense, and defense in turn attack.
To want to pin him down in a single one-sided mood,
confusing those brilliant, volcano-like eruptions
of idea-spouting polemic of mood with conservative
jeremiads of political despondency, is unreasonable; for
though he is rich in moods, he is never despondent;
in his fundamental mood there is a seething spring of thousand-hued
moods of the idea: he is a man with more than
one face, a bifrontal figure.

There is no one among us who has expressed himself so strictly
conservatively and yet so extremely revolutionarily, no one who
has to such a degree defended and to such a degree attacked the
established order as S. Kierkegaard. What a contrast,
what a crying difference between his first and
his last utterance on the relation, mark you the fundamental relation,
between religion and state! That he was not
a political but a religious author is in this respect
beside the point.

In the above-cited article against Dr. Rudelbach he has
as good as nothing to object to the existing
church constitution; for the time being everything is good, very good,
as it is. It is not the protection of the state, not
the ecclesiastical arrangements, but only a little worldly
levity that is to blame for Christianity at times
% --- p. 20 ---
\setcounter{footnote}{0}%
\opage{20}sinking down to what Dr. R. calls \emph{Habit-Christianity.}
``I hate,'' says S. K., ``Habit-Christianity
in whatever form it shows itself.''
This ``in whatever form'' I could wish
particularly noted; for Habit-Christianity can indeed
have many forms. And if there were nothing else to
choose, if the choice were only between this kind of Habit-Christianity:
a worldly levity that lives on carelessly in
the conceit of being a Christian, without perhaps even
having any impression of Christianity, and the kind of
habit Christianity found among sectarians, the awakened,
the hyperorthodox, partisans when it comes to the crunch,
I choose unconditionally the first. The first kind has indeed
only frivolously and negatively taken Christianity in vain,
if one can even always dare to judge it
so strictly; the second has perhaps in spiritual
pride, in any case positively taken Christianity
in vain; the first one could almost be tempted to
smile at, because one hopes; the second can make one
shudder.'' Here there is no thought that there could be
anything amiss in the whole state-church system, not the remotest
hint that the clergy's salaries as royal
officials should stand in conflict with their calling
as servants of the Word. But how long did it last
before the leaf turned! Four years later we read in
``The Moment, No. 3, 6'' the following announcement. ``As long
as there exist in Denmark 1000 royal livings
for teachers of Christianity, the best possible has been done
to prevent Christianity. As long as there are 1000 royal
livings, there will always be a corresponding
number of people who intend in this way to
% --- p. 21 ---
\setcounter{footnote}{0}%
\opage{21}earn their bread. Among these there will be a few,
a very few, who perhaps did find a calling to proclaim Christianity.
But just at the moment when it was to
become really serious, when, trusting in God alone and at one's own risk, one
would have to take it upon oneself to appear as a teacher --- just then the
state opens to them the comfort of taking a royal
office, whereby these few are, Christianly, stunted. The
far greater number have no calling at all to
proclaim Christianity, but regard it quite simply as
a livelihood. In this way the state succeeds in filling the whole
country with corrupted Christianity, the very greatest difficulty
for the placing of true Christianity, far
greater than complete paganism.''

No one who really wants to fix his attention on
these highly dissimilar utterances can, however much respect
and sympathy he otherwise feels for the exceptionally
gifted author, conceal from himself that here
is a contradiction. For here the question is not of different
better or worse conditions, since such a change concerns
only reality; no, here the question is of the idea as
idea, of a principle valid for all conditions. If it
could in the year 1851 accord with the essence of Christianity
that its duly appointed teachers be paid by the state,
then the same must also hold for the year
1855, and conversely. Either --- or: either Christianity must
in the course of four years have changed its principle, or
else the author must in the interval have wholly
altered his judgment about the essence of Christianity; since now
the first is unthinkable, and the last improbable,
the question forces itself involuntarily upon us: how
is the contradiction to be resolved?

% --- p. 22 ---
\setcounter{footnote}{0}%
\opage{22}To the new edition of ``Practice in Christianity,'' which
appeared in the midst of the strife, the author appended (``The Fatherland''
May 16, 1855) indeed some remarks clarifying
in this respect. ``This book, I have,''
says he, ``let go out in a wholly \emph{unaltered} edition,
because I regard it as a historical document.
If it were now to have come out, now that I have convinced
myself that the established order, Christianly, is untenable,
it would have been altered\dots{}. My earlier
thought was: if the established order is to be defended,
this is the only way: poetically (therefore by a
pseudonym) to place the judgment upon it, by thus
drawing on ``grace'' in the second power, so that it became Christianity,
not merely finding in ``grace'' forgiveness for
the past, but finding in grace a kind of indulgence
for the actual following of Christ and the actual
exertion of being a Christian. In that way
truth does after all enter into the established order, it defends
itself by judging itself, it admits the Christian
requirement'' etc. But that this clarification does not
make the matter clear, anyone can tell himself. It % print: Sageu (as printed)
is indeed a veritable trade in indulgences that the single individual here, if
in perfect disinterestedness as regards advantage
--- underhandedly and in all quiet --- has wanted to carry on
with the established order.

Of ``a worldly levity which in the year 1851
lives on in the conceit of being a Christian without perhaps
having any impression of Christianity,'' he says
that it has only frivolously and negatively taken Christianity
in vain and is --- ``if one can even always dare to
judge it so strictly'' --- to be smiled at, ``because one
% --- p. 23 ---
\setcounter{footnote}{0}%
\opage{23}hopes.'' In this there hides, unmistakably, the ulterior thought
that if the established order would only admit its
``worldly levity'' and concede ``the requirement,'' then
it should --- even without its confession of its
distance being able (see the cited article in ``The Fatherland'') to be ``called
a striving in the direction of coming closer to the requirement''
--- still have leave to call itself a Christian
established order and yet be exempted from ``the actual
following of Christ.'' A strange traffic between idea
and reality!

In the year 1855, when the author of ``Practice in Christianity''
has convinced himself that the established order, Christianly,
is untenable, the time of grace is over. The established order
is pronounced under the ban, the whole of Christendom is laid under
interdict. Still in the year 1854 the established order
could have saved itself by ``drawing on grace in the second power'';
but in 1855, when the Moment came, the Moment No. 1-2-3,
it was over with the possible indulgence, over, over with
drawing on grace. Reality might then fold
its hands, join in with a credo, creep on its knees, clothe
itself in sackcloth and ashes, in short, behave as it would, it
was over. In the time of Innocent the Third, when France
was to be laid under interdict, the legate appeared,
clad in a sky-blue stole, as on the day of the Savior's
passion, had the images of the saints shrouded in crepe and
stones thrown from the pulpit down among the congregation,
as a testimony that thus the Lord of Heaven and
Earth had now rejected the whole people. It was a
gripping sight, a spine-chilling scene. But for
serious thinkers it was still more spine-chilling
to see that single individual with the New Testament
% --- p. 24 ---
\setcounter{footnote}{0}%
\opage{24}in his hand, a penitent who was going in to judgment,
a dying man, already clad in the shroud, throw stones at
the old church and break the staff over the whole existing
``Christianity of Christendom.'' It happened wondrously
suddenly and yet with an eternity's earnestness; it was like
a rash act and yet so well considered, so long, so
all-sidedly prepared. A fearful break between idea and reality!

Such leaps from one extreme to the other
can be explained only by a violent tension, an inner
strife in the principle itself. At first it looks as if the Idea,
as regards everything external, were perfectly satisfied with
the given reality. The Idea seeks the single individual. The
single individual has enough in himself, for he has enough in the Idea.
The only thing he needs for inwardness's sake is
a few highly beneficial exercises in patience, self-denial,
forbearance, longsuffering, peaceableness, etc.,
but in this respect the existing reality will
at any time be able to be of service to him; for it is
just the good thing about reality: it seldom lets there
be a lack of some fine occasion to practice
those virtues, invaluable for the single individual's freedom and peace.
In this way the Idea toward the external surroundings
is mildness, pliancy, condescension itself. No
abolition for the single individual's sake, no inconvenience
with the external conditions; all is good, very good,
excellently good, as it is. But what does this mean?
It means that the Idea in absolute inwardness is infinitely
sure of itself, but it by no means means that
it is on good terms with reality. Far from it!
it means that the single individual with concealed contempt for all
% --- p. 25 ---
\setcounter{footnote}{0}%
\opage{25}that is external looks smilingly down upon the existing order, in the
tacit presupposition that the existing order with silent admiration
looks up to him, and humbly acknowledges the Idea's superiority.
Since now reality unfortunately is not always
so sensitive, and the existing order prefers to follow its
own head, it can happen at the slightest unforeseen occasion
that a violent war breaks out. For the Idea
makes great demands, the Unconditioned will rule. As
graciously mild as the gracious Idea can otherwise be,
so terribly severe in turn is the wrathful Godhead. Where
everything before suited the Idea, now nothing suits: reality
is condemned, a court-martial is held over the existing order.
The struggle is waged on unequal terms. The existing order
has physical power as its defense; the single individual,
the Idea's servant, fights only with spiritual weapons. It
is the Idea's glorification, its highest revelation, that
its true combatants can conquer inwardly by
losing outwardly, that more and more single individuals can become martyrs.
The struggle can be continued so long, so many victims
can fall, that the clumsy, brutal reality
at last for its own sake is compelled to crawl to
the Cross and make concessions to the Idea, on which condition
there can then be concluded, if not a real peace,
yet an armistice. But whether the condition
is otherwise peaceful or unpeaceful, the secret has
come to light. The relation between idea and
reality is namely so determined that the Unconditioned
must lie in an eternal strife against its conditions.
Within the secret there sits a contradiction.
The contradiction is double-edged: it cuts through reality,
and it wounds the Idea.

% --- p. 26 ---
\setcounter{footnote}{0}%
\opage{26}From the side of reality the contradiction is this, that reality
with its web of conditions incessantly hungers
for the Idea and yet cannot digest the Idea, since the
finite conditions can never be reconciled with the Unconditioned.
In the Idea itself the contradiction is this, that it wants to
be realized, that it inwardly desires reality without
yet being able to endure reality, since the Unconditioned
is constantly repelled by the conditions. If a
true mutual reconciliation between idea and reality is to be
possible, then the principle itself must be changed.

It is a fundamental unity of the Unconditioned with its
conditions on which everything essentially depends. There must
be in the Idea determinations of the conditioned that fit
external reality, and there must be given in reality's own
essence determinations of the Unconditioned that at least
approximately --- that is, under a continued development
--- satisfy the Idea. The question now is how
matters stand with these reciprocal determinations.

As for S. K., it has hitherto been and
still is our full conviction that he with his sharp
eye saw very well the reconciling intermediate determinations,
but that in accordance with the tactics
corresponding to his peculiar task he considered himself obliged to keep them
in the background. He had to hold strictly to the Unconditioned
and see to it that he worked by a sudden blow, if
he was to succeed at all in rousing the age out of
its spiritual torpor and in placing a decisive point. ``For,'' says
he, (The Moment No. 1) ``to place a decisive point
--- and this is the task --- cannot be done
in the same way as everything else; and when moreover
the misfortune of the age is this ``to a certain degree,'' to a certain
% --- p. 27 ---
\setcounter{footnote}{0}%
\opage{27}degree of going in upon everything, when this is just the sickness,
then above all care must be taken that, so far as possible, it does not happen that it too
to a certain degree goes in upon this, whereby all is lost. No, something decisive
is delivered otherwise than anything else. Like the beast of prey's
spring upon its quarry, like the eagle's stoop in its fall, so
something decisive is delivered: suddenly and gathered into
One (intensively). And as the beast of prey unites cunning with
strength, first cunningly holds itself quite still --- as still
as no tame animal can be; and then gathers itself wholly
into a single spring or stroke --- as no
tame animal can gather itself for or rise into the spring: so
something decisive is delivered. First still --- so still is
the weather never on a still day, so still it is only before
thunder; and then it breaks loose''.

The basic form that alone is capable of taking up
the intermediate determinations of the conditioned and the unconditioned
and of reconciling the Idea with reality is, with its
rich content, more than a form, it is itself an idea,
a true idea of reality, namely the idea of society. It
is the idea of society that S. K. everywhere puts in shadow,
so that a so much stronger light may fall upon his single individual.
``The single individual,'' he says, ``is the category through
which, in a religious respect, the age, the generation, history
are to. And he who stood at Thermopylae was
not so secure as I, who, in order at least
to get attention drawn to it, have stood at this narrows `the single individual.'
He was namely to keep the throngs
from pressing through the narrows; if they pressed through,
then he had lost. My task is at least
far less exposing me to the danger of
% --- p. 28 ---
\setcounter{footnote}{0}%
\opage{28}being trampled down, since it was that of a lowly servant\dots{} if
possible to cause, to invite, to move the many
to press through this narrows `the single individual,'
through which, mind you, no one presses except
by becoming the single individual; the opposite is indeed a categorical
impossibility.'' Here, then, it is assumed that the narrows
``the single individual'' is quite otherwise secure than the narrows
at Thermopylae; but about this one may be permitted
to harbor a doubt.

The narrows at Thermopylae had, after all, always, whether
the throngs pressed through or not, its definite
extent, its rock-firm boundaries; but the narrows
``the single individual'' has, so far as we can see, no
rock-firm boundaries: it constantly widens as
the throngs press through, even to such a degree that it
finally must hold them all. For when all become the single individuals,
then the single individuals must necessarily again become all,
but to become all is indeed to become the crowd: by going
out of the crowd --- into the narrows, the throngs go
thus precisely into the crowd --- out of the narrows.
In order to enter into the truth the single individuals without exception must separate
themselves from the crowd, for the crowd is untruth;
but when now each single individual, after they all have
separated themselves from the crowd, just for that reason comes
to stand in the midst of the crowd: where does he then find himself, in
the truth or in untruth? If we say in untruth,
since he is one of the crowd, and the crowd is untruth,
then a mistake has been committed here in the narrows, here is,
frankly speaking, a disaster come to pass. When fire breaks out in a
large building, where many are assembled, and the building's
only exit is a narrows, it easily happens that
% --- p. 29 ---
\setcounter{footnote}{0}%
\opage{29}all burn inside, because all want out at once: by
going singly, one after one, at least
some could have been saved. Roughly so it stands
with the narrows ``the single individual,'' when it is to hold
unconditionally that crowd is untruth. Only by being
the exception, the one distinguished above the others, can the
single individual categorically separate himself from the crowd. Every
generation, every throng one must therefore, when the single individual
is to come forward, cry out to: One at a time, one in each generation,
only one, a single one! for one more and one more is already
crowd.

It is evident that the crowd cannot possibly form
what S. K. calls a \emph{qualitative} opposition to the
single individual. The crowd is formed of single individuals; as the single individuals
are, so is the crowd. If the crowd consists of true
single individuals, then there is truth in the crowd; if the crowd is
untruth, then the ground is that it consists of
untrue single individuals. The crowd as crowd is in and of itself
neither good nor evil, neither true nor untrue. In order
to grasp the crowd as the existing untruth, one must
first in one's concept exclude it from the truth, by
excluding it from the Idea. The idealess crowd --- yes,
that is indisputably true --- the idealess crowd is, as the case may be,
mass, rabble, public, the numerical,
the untrue; but it is then not the untrue
because it is crowd, it is so only because it is excluded
from the truth. In a corresponding way one can
indeed also begin by excluding the single individual from
the truth, regarding the single individuals outside the Idea; and it
then follows quite of itself that the idealess, that is to say
untrue, single individuals are untrue. But the single individuals are then
% --- p. 30 ---
\setcounter{footnote}{0}%
\opage{30}untrue, not because they are single individuals, but because by
the concept set up they are excluded from the truth.

If any true opposition is to be formed between
the crowd and the single individual, then both opposed
terms must fall within the Idea; but the crowd, when it is referred
to its idea, is existing society. Thus
we come to the opposition between the single individuals and
society; crowd is reality, society idea: we
must then lay particular weight on the question: how
things stand in real society with idea
and reality.

Society has its subsistence in the organized
crowd. When now the Idea is to develop and realize
itself in society, which way does it then go? Does it begin
partly with the single individuals and work its way first little by little
through the whole, or does it begin straightway
from the whole, in order gradually to kindle its light in
the single individuals, or does it begin at one stroke both from
the whole and from the single individuals? To these questions, so highly
important for the idea of freedom, S. Kierkegaard
gives us nowhere any definite answer. He pushes the investigation
from him, presupposes without further ado society's
subsistence, polemicizes against the crowd and holds to the
single individual. But that the investigation, precisely for the single individual's
sake, is of the utmost importance will become evident when
we cast a glance at the position of the single individuals in the diverse,
narrower and wider circles of society, in the congregation,
in the state, in the people.

If we, to begin with the religious, set the single individual
against the congregation, then we meet at once the fact
that so long as the truth existed exclusively as
% --- p. 31 ---
\setcounter{footnote}{0}%
\opage{31}the prototype himself, the One, the unconditionally single individual,
the congregation could not be founded, and when the congregation was
founded, no single individual could step forth independently against the congregation's
spirit. But in zeal to defend the narrows
``the single individual'' S. K. goes in this connection so far
as even to assert that ``Christianity has really
not come into the world,'' since not even the apostles
have properly been able to keep order at the entrance. ``For
one thing is,'' he says (The Moment No. 5), ``to work in such a way
for dissemination that one uninterruptedly early and late proclaims
the doctrine to all, another thing is to be too quick
to permit people thus in hundreds, in thousands to assume the
name of Christians\dots{}. Had a multitude of people
let themselves be seized by Christ's speech, he would
assuredly not at once have permitted these thousands to call themselves Christ's
disciples. No, he held back more strongly. He won
therefore also in 3½ years only 11 --- whereas an apostle
in \emph{one} day, perhaps in \emph{one} hour, wins 3000 disciples of Christ.
Either the disciple is here greater than the master, or the truth
is that the apostle is a little too quick to strike,''
etc.

If this is taken according to the letter, then the idea of the congregation
is without further ado to be given up. The idea of the congregation
namely presupposes that the spirit at one and the same time begins its
activity from society and from the single individuals.
The apostle speaks as the single individual to the single individuals in the crowd;
but he speaks from the spirit. The spirit from which he
speaks is the religious spirit of society, the spirit that
will not work singly upon the scattered single individuals, but
unite them in society, in order in and through society to
work in every single individual thereby that they could work upon
% --- p. 32 ---
\setcounter{footnote}{0}%
\opage{32}one another reciprocally and under reciprocal interaction
develop into true single individuals. Therefore he hastens to
get the congregation founded. Had the apostle waited.
until the single individuals had become more prepared, so that
the truth in the congregation could at last become a plain and
simple result of the truth in the single individuals formed
outside the congregation, then he would never have come to found
any congregation. When, then, should he have founded it?
Before the single individuals had yet become true single individuals? That
was too early! After they had become true single individuals?
That was too late! It is another matter if the apostle
had let it stay with baptizing the many thousands
and said to the crowd: now you are baptized, now you are the congregation,
now nothing more is needed; but that is not how it
went. It is so far from the case that the crowd as
crowd comes to rule in the congregation, that the congregation's
spirit precisely sets its single individuals, its chosen,
the great task of setting the crowd right.
of organizing, disciplining, bringing it up, so that the many, each for himself,
must recognize what it means to be before God the
single individual.

If we, as regards the political, set the single individual
against the state, a similar result emerges. The mass,
the idealess crowd can naturally not found any
state; but the single individual who despises the crowd cannot
either. The idea of the state has two equally original
points of departure. It already stirs powerfully in the crowd,
in that it begins to clarify itself for the individual
leaders. In order to organize the masses and work against the
public, the merely numerical, it is not enough ``to cause,
to invite, to move the many to press through
% --- p. 33 ---
\setcounter{footnote}{0}%
\opage{33}the narrows the single individual,'' it is at the same time necessary
that one uninterruptedly early and late inculcate
the validity of the idea of society, the authority of the government, the
rights of the citizen and the duties of the citizen. By separating
the consciousness of civic duties from religious
inwardness and impressing upon the single individuals contempt for the tasks of society
under the pretense that these are only
worldly external forms, one does the state a poor
service. What does it help to bow to the authorities
and speak piously, submissively, of the magistracy being
appointed by God, and of the administration ought to be left to ``those duly
appointed and called thereto,'' when in the same
moment one undermines the foundation of all rational government?
Who is it that is to govern? the single individuals or
not the single individuals? If it is not the single individuals, then
it must indeed be the crowd, the public, the numerical.
But if it is only ``the leaders of the crowd'' who make up
the government, then it sounds strange indeed that such a
government ought not essentially to seek its strength in
``worldly prudence,'' but in ``an inward transformation in the
direction of becoming steadfast in fearing God.'' If the government is to
become steadfast in fearing God, then it must necessarily
consist of single individuals duly appointed thereto,
of such as in religious inwardness have learned
what it is to press through ``the narrows.'' In these
true single individuals the religious cannot possibly be
indifferent to the political; in order to gain strength to
govern, they must become steadfast in fearing God, but the
power of inwardness they win by fearing God shall in turn
show itself in the power and authority with which
they govern: a more inward interaction between politics
% --- p. 34 ---
\setcounter{footnote}{0}%
\opage{34}and ``fear of God'' cannot be thought. What holds
of some true single individuals holds of all; they are indeed
all to go through the narrows. In the capacity of citizens
society's single individuals now have the finest
example in the government. Just as the ``God-fearing'' government
for its part must in no way regard the outer
social forms as indifferent, notwithstanding that they are worldly,
so the God-fearing citizens for their part dare not
either. If it is a matter of conscience for a government composed of true
single individuals to further the development of the idea of society
and to care for the welfare of the state by watching over
the rights of the citizens, then it must also be a matter of conscience for all individual citizens to live themselves into the
idea of society and punctually fulfil the civic duties.

Although the religious and the political idea must indeed
be opposed to one another, as inner freedom is opposed
to outer, they nevertheless have this in common, that both are
ideas of society, both laid out to engage with
the crowd, take the conditions up into the unconditioned and transform
reality. The role the crowd plays in the
determination of the relation between idea and reality
shows itself especially prominently when the talk is of society
in its widest extent, of the \emph{people.}

That the \emph{people} in a state-society must, when we look at
reality, exist as crowd cannot be denied;
but just for that reason it becomes, for the sake of the Idea,
of the greatest importance to make a distinction between the people and
the crowd. The people is the crowd insofar as it is the sum
of the population's individuals, and yet the crowd is not yet
the people; only by being referred to its original idea,
the idea of the fundamentally human in its natural-force
% --- p. 35 ---
\setcounter{footnote}{0}%
\opage{35}peculiarity (cf. ``On Obstacles to and Conditions for Spiritual Life in the Present Age,'' VIII), can the crowd be conceived
as a people. Were it really impossible to make the distinction between
the people and the crowd valid in a penetrating,
clear manner, evident to sound reason,
then freedom of the people and freedom of the crowd would be synonymous
words. But if freedom of the people without more falls together
with freedom of the crowd, so that a free constitution is to
be one that sets up the conditions for the crowd lawfully
to get its will, even if this will be in itself ever so
unreasonable, put through: then it is not to be seen how any reasonable human being
could place either trust or hope in such a
freedom. It is only the opponents of the freedom of the people who,
whether they otherwise despise or idolize the crowd,
make the people one with the crowd.

In what sense the people as people must be said to
be raised above its own crowd and the true will of the people
exalted above the accidental will of the majority can easily be explained by
some examples. Suppose that the majority in a population, which is
hardly likely, but yet not unthinkable, to such a degree
surrendered itself to faith in eternal peace that it regarded
all defense as a superfluity, and on that occasion
got put through that all items of expenditure for army and
navy were struck out: then it would still not be the people
that harbored such a faith, it was only the crowd.
Or suppose that a vote once with a great majority
brought about permission that men and women ---
for the sake of personal liberties --- henceforth might
dress alike, so that woman, in order to come to her
natural right, might even dare to put on a false beard, so
% --- p. 36 ---
\setcounter{footnote}{0}%
\opage{36}then it was by no means the people that came to speak through such a
vote, no, it was only the crowd. Or
suppose that a liberal-minded majority, in recognition that marriage
is as much a civil as a natural institution,
got put through at a twenty-seventh reading of the marriage law
the change that it should henceforth
be permitted a man to have several wives,
provided only that the number of wives came to stand in strictly regulated
proportion to the income tax, then a progress of
this nature would still not be attributable to the people, but to
the crowd. What is it now that in the examples cited
makes it so easy to distinguish the crowd's whims from the people's
thoughts? It is precisely the in itself absurd,
unreasonable, for the whole harmful; for the people as people
can will only the reasonable, what accords with life's basic demands,
what is beneficial to the whole of society.

If all matters important for society were
of such a nature that nothing more were needed than sound
sense, natural understanding and a good will in order at once to
see what it is on which it each time depends,
and how much in critical moments may be at stake,
then it would still be an easy matter to clarify the people's
view and set the crowd right. But the difficulty is
that precisely the most significant tasks, the great questions
that reach deepest into the life of society, the matters that are popular
in the most proper sense, belong to those
whose significance is most difficult to grasp and fully to recognize
by the whole people, because the crowd, as the saying goes, cannot
keep up. If this hitherto insurmountable obstacle to true popular cultivation
is at last to be overcome, it can
happen only on the condition that the popular
% --- p. 37 ---
\setcounter{footnote}{0}%
\opage{37}questions at one and the same time are treated both theoretically
as questions of the idea and practically as questions of reality:
idea and reality must then not be separated, but
held against each other, so that in clear, simple outlines it can
be seen how in every matter concerned they illuminate
each other reciprocally.

The question of the religion of the people is, as surely everyone
will concede, both a question of the idea and a question of reality.
The common opinion, in itself quite
correct, that a people cannot possibly do without a positive
religion, is too vague and weak to bear the reality's
heavy, deep-cutting debates concerning
the position of the people's church toward the state. The first thing needed
is that the people in its entirety come to clarity about the essence of religion
and see its significance for human life.
Only gradually as a living cognition
of what religion is, in what inner essential relation
it stands to every independent people's morals and way of thinking,
to the progress of culture and civilization in all
directions, spreads more and more, will the people appear
as a people and lay the weight of its mood and voice
on the scales of the debates.

What distortions of words and misunderstandings can be
set in circulation when great questions of the idea are without more
drawn down into finitude and furnish slogans for the benefit
of the parties, is seen from the present age's peculiar dispute about
the worth of intelligence. Outside the Idea such a dispute will
never end. It is evident that those who cast a shadow
upon intelligence connect with the word an entirely different
conception than those who defend it. Intelligence is
insight, insight requires enlightenment. Without enlightenment
% --- p. 38 ---
\setcounter{footnote}{0}%
\opage{38}no freedom, without popular enlightenment no freedom of the people.
What is freedom without enlightenment, a free people without
intelligence? A day without light, a sword without blade, ---
thoughtless contradictions! Against imagined cultivation, the
empty enlightenment, the puffed-up, unfruitful intelligence
objection ought to be made; but well-founded objections
against false intelligence indeed contain precisely a defense
of the true.

On the surface of reality it may indeed in the eyes of the crowd
sometimes look as if art, poetry and
science were only luxury articles for the enjoyment of the
higher classes of society; but a people that wants to understand
itself will also understand that the clear forms of the spirit
belong essentially to the free, many-sided development of the life of the people,
are true necessities, although they cannot, like
``\emph{Bread and Iron},'' straightforwardly be reckoned among the first
necessities. The basic question is how the
higher cultivation as a whole relates to the lower;
for that the lower cannot raise itself without help from the
higher is a fact. Ask the worker, he
already knows that one can work properly only
with the hand when one at the same time works in the spirit, that
he can use his strengths fully only when he has
learned to use his thoughts.

A telling proof that enlightenment and freedom go hand
in hand is the unmistakable need, the drive and striving
that everywhere stirs in the population after knowledge
and cultivation: the masses gather where they sense a thought;
they listen to the word, they glimpse the light, they press forward
to greet the day of freedom. This striving, this need
is a testimony that the people as people has a deep
% --- p. 39 ---
\setcounter{footnote}{0}%
\opage{39}sense of its future's task, that it will not
settle down at rest with outward freedom, but stride
further forward and know how to win by struggle the inward.
Only by the union of outward freedom with the inward
does popular freedom become complete, conditioned in
reality, unconditioned in the Idea: the freedom of the people is, like
the people itself, a becoming, growing unity of idea and
reality.
\signature{R. \emph{Nielsen}.}

\end{document}
